% Preamble of the article.  Chapter files use only these packages and macros.
\documentclass[11pt,a4paper]{report}
\usepackage[T1]{fontenc}
\usepackage[utf8]{inputenc}
\usepackage{lmodern}
\usepackage[margin=2.7cm]{geometry}
\usepackage{amsmath,amssymb,amsthm,mathtools}
\usepackage{mathrsfs}
\providecommand{\llbracket}{[\![}
\providecommand{\rrbracket}{]\!]}
\usepackage{booktabs,longtable,array}
\usepackage{enumitem}
\usepackage{xcolor}
\usepackage[hidelinks]{hyperref}
\usepackage[nameinlink,capitalise]{cleveref}

\theoremstyle{plain}
\newtheorem{theorem}{Theorem}[chapter]
\newtheorem{lemma}[theorem]{Lemma}
\newtheorem{proposition}[theorem]{Proposition}
\newtheorem{corollary}[theorem]{Corollary}
\theoremstyle{definition}
\newtheorem{definition}[theorem]{Definition}
\newtheorem{example}[theorem]{Example}
\newtheorem{question}[theorem]{Question}
\theoremstyle{remark}
\newtheorem{remark}[theorem]{Remark}

% --- names ---
\newcommand{\iPRA}{\mathsf{iPRA}}            % intuitionistic primitive recursive arithmetic
\newcommand{\Teq}{\mathsf{T}_{\mathrm{eq}}}  % the ordinary scalar intermediary
\newcommand{\Nat}{\mathsf{N}}                % the native scalar core (eight-rule kernel)
\newcommand{\pkg}[1]{\texttt{#1}}            % a research package, e.g. \pkg{witness\_rule}
\newcommand{\Lean}[1]{\texttt{#1}}           % a Lean declaration or module
% A rule name, e.g. \rulename{Use}: every letter in small capitals, so that
% all-capital names such as QE, DE, DI stand apart from ordinary capitals.
\newcommand{\rulename}[1]{\textsc{\MakeLowercase{#1}}}
% --- terms and formulas ---
\newcommand{\It}[4]{I(#1,#2;\,#3.\,#4)}      % the iterator I(n,a;u.h): \It{n}{a}{u}{h}
\newcommand{\Suc}{S}                         % successor
\newcommand{\sel}{\operatorname{sel}}        % selector sel(c,a,b) = I(c,a;u.b), u not in b
\newcommand{\D}[2]{D_{#1,#2}}                % the paired binder D_{x,y}: \D{x}{y}B
\newcommand{\Q}[1]{Q_{#1}}                   % the closure binder Q_y: \Q{y}B
\newcommand{\nf}{\operatorname{nf}}          % normal form
\newcommand{\FV}{\operatorname{FV}}
\newcommand{\proves}{\vdash}
\newcommand{\subst}[2]{[#1/#2]}              % substitution t[s/x]: t\subst{s}{x}
\newcommand{\defeq}{\mathrel{:=}}
\newcommand{\iffdef}{\mathrel{:\Leftrightarrow}}
\newcommand{\infrule}[2]{\dfrac{\;#1\;}{\;#2\;}} % an inference: \infrule{premises}{conclusion}
